\section{Model}
\label{sec:model}

\subsection{The unit transfer raiding game}

The analysis begins from a repeated raiding game whose only purpose is to make the opening reciprocal order exact. Time is discrete and indexed by $t = 0,1,2,\dots$. Agents are indexed by $i \in \{1,\dots,N\}$ with $N \ge 3$, and wealth at the beginning of period $t$ is the publicly observed vector $w_t = (w_{1t},\dots,w_{Nt})$. In the benchmark all agents begin with common wealth $\bar{w} > \tau$, where $\tau > 0$ is the unit amount that can be stolen from an unoccupied house in one period. The richer proportional environment appears only after the benchmark and keeps the same state variable. The notation used throughout the paper is collected in Table \ref{tab:notation}; the same symbols are used in the theory, the simulations, and the appendix proofs.

\begin{table}[tbp]
    \centering
    \caption{Notation. The table collects the state variables, contract terms, and regime objects used throughout the theory and simulation sections.}
    \label{tab:notation}
    \begin{tabular}{p{0.17\textwidth}p{0.73\textwidth}}
\toprule
Symbol & Meaning \\
\midrule
$i \in \{1,\dots,N\}$ & Agent index. \\
$t = 0,1,2,\dots$ & Time index. \\
$w_{it}$ & Wealth of agent $i$ at the beginning of period $t$. \\
$\sigma$ & Baseline theft permutation defining the reciprocal cycle. \\
$h$ & Exogenously honest agent in the perturbation analysis. \\
$\tau$ & Unit loot in the benchmark stage game, normalized to $\tau = 1$. \\
$\gamma$ & Proportional theft intensity in the dynamic extension, with $\gamma \in (0,1)$. \\
$c_T$ & Resource cost of personal theft effort. \\
$\rho$ & Exposure coefficient capturing the expected fraction of own wealth lost when one leaves home. \\
$\phi_t$ & Targeting advantage of delegated theft relative to freelance theft. \\
$s_{ijt}, \alpha_{ijt}$ & Fixed payment and loot share in a patron-thief contract between patron $i$ and thief $j$. \\
$o_{jt}$ & Outside option of poor agent $j$ if he declines a patron's contract. \\
$k$ & Contracting cost borne by the patron. \\
$g$ & Guard wage. \\
$q_t$ & Probability that guarding blocks a raid on protected wealth. \\
$R_t, E_t, P_t$ & Sets of rich, elite, and poor agents at date $t$. \\
\bottomrule
\end{tabular}

\end{table}

In each benchmark period, agent $i$ chooses one action from the set
\[
A_i = \{H\} \cup \left(\{1,\dots,N\}\setminus\{i\}\right),
\]
where $H$ means staying home and $j \neq i$ means leaving home and attempting to raid household $j$. Occupancy determines vulnerability. If $a_{jt}=H$, household $j$ is occupied and no raid on $j$ succeeds in period $t$. If $a_{jt}\neq H$, household $j$ is unoccupied and can be raided. This is the only occupancy rule needed for the benchmark, and it makes the first honest night deterministic once one household stays home permanently.

Target congestion is explicit. For any household $j$, let
\[
M_{jt} = \{i \neq j : a_{it}=j\},
\qquad
m_{jt} = |M_{jt}|.
\]
If $a_{jt}=H$, every attempted raid on $j$ fails. If $a_{jt}\neq H$ and $m_{jt}\ge 1$, at most one unit $\tau$ can be extracted from $j$ during period $t$, and that unit is divided in expected value equally across the raiders in $M_{jt}$. The benchmark therefore treats loot as rival and rules out the possibility that several agents can each take the same unit independently.

Define the realized loss of household $i$ and the expected loot of agent $i$ by
\begin{equation}
\ell_{it}(a_t)
=
\tau \mathbf{1}\{a_{it}\neq H\}\mathbf{1}\{m_{it}\ge 1\},
\qquad
g_{it}(a_t)
=
\sum_{j\neq i}\mathbf{1}\{a_{it}=j\}\mathbf{1}\{a_{jt}\neq H\}\frac{\tau}{m_{jt}}.
\label{eq:loss-gain}
\end{equation}
The benchmark abstracts from direct theft cost so that the reciprocal accounting is exact. End-of-period wealth is therefore
\begin{equation}
w_{i,t+1} = w_{it} - \ell_{it}(a_t) + g_{it}(a_t).
\label{eq:benchmark-wealth-update}
\end{equation}

\subsection{Balanced reciprocity}

\begin{definition}[Balanced theft cycle]
\label{def:balanced-cycle}
A balanced theft cycle is a permutation $\sigma$ of $\{1,\dots,N\}$ with no fixed points such that each agent $i$ chooses $a_i=\sigma(i)$.
\end{definition}

Under a balanced theft cycle every household is unoccupied, every household is targeted by exactly one raider, and no target congestion appears because $m_j=1$ for every $j$. The basic accounting is then immediate but worth stating explicitly because the rest of the paper measures departures from this object.

\begin{lemma}[Balanced accounting]
\label{lem:balanced-accounting}
If all agents begin a period with common wealth $\bar{w}$ and play a balanced theft cycle, then every agent ends the period with wealth $\bar{w}$ and aggregate wealth is preserved.
\end{lemma}

\begin{proof}
See Appendix \ref{app:proofs}. Each agent loses exactly $\tau$ and gains exactly $\tau$, so wealth is unchanged agent by agent and therefore in aggregate.
\end{proof}

\begin{proposition}[Balanced reciprocity as a weak equilibrium]
\label{prop:balanced-cycle}
In the one-shot unit transfer stage game, any balanced theft cycle is a weak Nash equilibrium.
\end{proposition}

\begin{proof}
See Appendix \ref{app:proofs}. Deviating to staying home preserves the deviator's wealth but does not raise it, because the avoided loss is exactly offset by forgone loot. Deviating to a different target creates congestion, so expected loot falls below $\tau$ while the incoming loss remains. Hence no unilateral deviation is strictly profitable.
\end{proof}

Only a weak claim is needed at this stage. The purpose of the benchmark is not to establish efficiency or to characterize the full set of equilibria. It is to show that a reciprocal predatory order can be a coherent strategic object once actions and congestion are written down exactly. The resulting graph is a directed cycle, and that graph theoretic resemblance to top trading cycles remains useful when one link later breaks. The resemblance goes no further. Transfers here are coercive, occupancy changes the opportunity set of others, and rivalry in extraction creates externalities that have no analogue in the voluntary exchange environment of \citet{shapley_scarf_1974}.

\subsection{A permanent honest household and stationary continuation}

The repeated benchmark keeps the same stage game and adds two restrictions immediately. The horizon is infinite with common discount factor $\delta \in (0,1)$, and strategies are restricted to be stationary and Markov in current wealth. A strategy for agent $i$ is therefore a map from the current state $w_t$ into an action in $A_i$. This is the equilibrium class studied in the theorem below. The argument does not characterize the full set of repeated outcomes that could be sustained by history-dependent punishments.

Fix one agent $h$ who is permanently honest and therefore chooses $H$ in every period. Let $\sigma$ denote the pre-perturbation balanced theft cycle, let $p = \sigma^{-1}(h)$ be the predecessor who planned to rob $h$, and let $q = \sigma(h)$ be the successor whom $h$ would have robbed under universal participation. If every other agent continues the old cycle on the first perturbed date, then the occupancy rule makes the first wealth change deterministic:
\begin{equation}
w_{h1} = \bar{w}, \qquad w_{p1} = \bar{w} - \tau, \qquad w_{q1} = \bar{w} + \tau,
\label{eq:first-night-accounting}
\end{equation}
and all other agents remain at $\bar{w}$.

\begin{lemma}[First night local asymmetry]
\label{lem:first-night}
Under the benchmark cycle with one permanently honest agent $h$, the first honest night creates a local three agent asymmetry given by \eqref{eq:first-night-accounting}. In particular, the successor $q$ receives a windfall and the predecessor $p$ suffers an immediate shortfall.
\end{lemma}

\begin{proof}
See Appendix \ref{app:proofs}. The honest agent neither steals nor is robbed, the predecessor fails in his attempted raid on an occupied house, and the successor is not robbed by the honest agent.
\end{proof}

The benchmark stage game is enough to show that exact reciprocity breaks. A withdrawal result needs one additional margin. Suppose that personal raiding in the repeated game carries direct effort cost $c_T \ge 0$ and expected exposure loss $\rho w_{it}$, with $\rho \in (0,1)$, because the raider leaves current wealth unattended while away from home. The expected one-period return from personal theft is then
\begin{equation}
\pi^{S}(w_{it}) = \tau - \rho w_{it} - c_T,
\label{eq:self-theft-threshold}
\end{equation}
while the normalized return from staying home is $\pi^{N}(w_{it}) = 0$. Hence there is a wealth threshold
\begin{equation}
\widehat{w} = \frac{\tau - c_T}{\rho}
\label{eq:threshold}
\end{equation}
such that an agent weakly prefers personal theft when $w_{it} \le \widehat{w}$ and strictly prefers staying home when $w_{it} > \widehat{w}$.

\begin{assumption}[Interior exposure threshold]
\label{ass:threshold}
The exposure threshold satisfies $\bar{w} < \widehat{w} < \bar{w} + \tau$.
\end{assumption}

The assumption places the symmetric benchmark wealth level below the exposure cutoff and the first-night windfall above it. Personal theft remains privately viable at $\bar{w}$ but ceases to be optimal after one unit of windfall. This is a narrow comparative static statement and not a general theory of criminal choice.

\begin{theorem}[No stationary one-honest continuation]
\label{thm:no-stationary-one-honest}
Under Assumption \ref{ass:threshold}, there is no stationary Markov equilibrium of the repeated benchmark game that preserves the original balanced theft pattern for all agents other than the permanently honest agent $h$.
\end{theorem}

\begin{proof}
See Appendix \ref{app:proofs}. By Lemma \ref{lem:first-night}, the successor $q$ enters period $t=1$ with wealth $w_{q1} = \bar{w} + \tau$. Assumption \ref{ass:threshold} implies $w_{q1} > \widehat{w}$, so \eqref{eq:self-theft-threshold} yields $\pi^{S}(w_{q1}) < 0 = \pi^{N}(w_{q1})$. Thus $q$ strictly prefers staying home at $t=1$ and no longer continues to steal. Since the original balanced pattern requires $q$ to keep stealing, that pattern cannot be stationary after the perturbation.
\end{proof}

The theorem is deliberately local. It does not solve the full repeated game with one honest agent. It says only that the original balanced cycle cannot survive as a stationary Markov continuation once the first honest night moves the successor across the exposure threshold. That is exactly the comparison later implemented in the deterministic benchmark simulation.

Chronology matters. The honest agent is not robbed immediately upon becoming honest because occupied houses are perfectly safe in the unit transfer benchmark. Only later, after the paper moves to the proportional environment with delegated predation and guarding, do occupied but unguarded houses become imperfectly safe. The deterministic first-night result and the later stochastic vulnerability therefore belong to different layers of the model and should not be conflated.

\subsection{Proportional-transfer extension and endogenous patronage}

Once the first honest shock has broken symmetry, the benchmark no longer suffices. The dynamic extension retains the same population, the same public wealth state, and the same logic of explicit target choice, but the theft technology changes. A successful raid on victim $j$ now transfers $\gamma w_{jt}$, with $\gamma \in (0,1)$. If several raiders converge on the same unprotected household, total extractable loot is still rival and concurrent raiders divide that amount rather than each taking it independently. This is the only theft technology used in the stratification, patronage, and guard regimes. Under unequal wealth, the relevant margin is no longer the gain or loss of one unit. A rich agent who leaves home exposes a large stock of wealth, while a poor agent exposes very little. The proportional specification expresses that asymmetry parsimoniously.

Let $m_t$ denote the expected gross loot from the best personal raid available to a rich agent at date $t$. For most of the analysis it is enough to treat
\begin{equation}
m_t = \gamma \bar{w}^{T}_t,
\label{eq:mt}
\end{equation}
where $\bar{w}^{T}_t$ is the expected wealth of the best available target class among nonelite households. If rich agent $i$ raids personally, his expected one-period net payoff is
\begin{equation}
\pi^{\mathrm{self}}_{it} = m_t - \rho w_{it} - c_T.
\label{eq:personal-payoff}
\end{equation}
If instead he stays home without contracting, the normalized payoff is $\pi^{\mathrm{stay}}_{it}=0$.

Patronage becomes meaningful only if delegation changes the feasible target set and the division of surplus. To capture the first margin, let $H_t \subseteq P_t$ denote the set of relatively wealthy but still poorly protected households among the poor. A patron can direct a hired thief toward this set more effectively than a freelancer can. Let $\phi_t \ge 1$ denote that targeting advantage, so delegated gross loot is
\begin{equation}
L_t = \phi_t \gamma \bar{w}^{H}_t,
\label{eq:delegated-loot}
\end{equation}
where $\bar{w}^{H}_t$ is the mean wealth of the targeted high value poor. The model thus keeps the direction of physical theft clear. In the patronage regime, poor thieves usually rob poor or near-poor households, not elite households. The net transfer is poor to rich because some of the loot is remitted upward to the patron.

To capture the second margin, let $o_{jt}$ denote the outside option of poor agent $j$ if he rejects an offer from a patron. The outside option is not the full delegated loot because a freelancer lacks the targeting advantage of the patron and competes with other poor thieves for a depleted prey base. The model writes
\begin{equation}
o_{jt} = \lambda_t \gamma \bar{w}^{P}_t - c_T,
\qquad \lambda_t \in (0,1),
\label{eq:outside-option}
\end{equation}
where $\bar{w}^{P}_t$ is mean poor wealth and $\lambda_t$ summarizes congestion and bargaining weakness in the freelance market. A contract between patron $i$ and thief $j$ has fixed payment $s_{ijt}$ and loot share $\alpha_{ijt} \in [0,1]$ for the thief, while the patron pays contracting cost $k>0$. The thief accepts if
\begin{equation}
s_{ijt} + \alpha_{ijt}L_t \ge o_{jt}.
\label{eq:IR}
\end{equation}

Let $\beta \in [0,1]$ denote patron bargaining power over the residual surplus. If a contract is formed, total surplus net of the outside option of the thief and the contracting cost of the patron equals
\begin{equation}
\Sigma_t = L_t - k - o_{jt}.
\label{eq:total-surplus}
\end{equation}
The payment rule is chosen so that the thief receives his outside option plus the share $(1-\beta)\Sigma_t$, leaving the patron with $\beta \Sigma_t$. Equivalently,
\begin{equation}
s_{ijt} + \alpha_{ijt}L_t = o_{jt} + (1-\beta)\Sigma_t
\label{eq:bargaining-rule}
\end{equation}
whenever $\Sigma_t \ge 0$.

\begin{lemma}[Contract feasibility]
\label{lem:contract-feasibility}
There exists a patron-thief contract satisfying the thief's participation constraint if and only if
\begin{equation}
L_t - o_{jt} > k.
\label{eq:feasibility}
\end{equation}
\end{lemma}

\begin{proof}
See Appendix \ref{app:proofs}. A contract is feasible exactly when gross delegated loot exceeds the reservation utility of the thief plus the contracting cost of the patron.
\end{proof}

Condition \eqref{eq:feasibility} is economically transparent. Patronage is not attractive merely because the patron is rich. It is attractive because delegated predation reaches a target class or success probability that the poor freelancer cannot reach alone, and because the outside option of poor workers is compressed by competition and depletion. If either margin disappears, the contract market collapses.

\begin{proposition}[Strict patron surplus and occupational choice]
\label{prop:patronage}
Suppose Lemma \ref{lem:contract-feasibility} holds and the patron captures fraction $\beta$ of the contract surplus. Then patron $i$'s expected payoff from hiring thief $j$ is
\begin{equation}
\pi^{\mathrm{hire}}_{ijt} = \beta\left(L_t - k - o_{jt}\right).
\label{eq:patron-payoff}
\end{equation}
The rich agent strictly prefers hiring thief $j$ to both abstention and personal theft whenever
\begin{equation}
\beta\left(L_t - k - o_{jt}\right) > \max\left\{0,\; m_t - \rho w_{it} - c_T\right\}.
\label{eq:patron-dominance}
\end{equation}
\end{proposition}

\begin{proof}
See Appendix \ref{app:proofs}. Equation \eqref{eq:patron-payoff} follows directly from the bargaining rule. Hiring dominates abstention because the right-hand side of \eqref{eq:patron-dominance} exceeds zero, and it dominates personal theft because the expected surplus from delegation exceeds the net return available to the rich agent who exposes wealth at home.
\end{proof}

Surplus arises because delegation is more productive than freelance theft and because the reservation utility of the worker remains below delegated gross returns. The wedge can come from better targeting, better information, bargaining power, or labour-market slack. The model requires only one such margin to survive.

Poor agents become hired thieves because they remain below the wealth threshold at which exposure makes personal theft unattractive and because the wage or commission implied by \eqref{eq:bargaining-rule} exceeds their shrinking outside option. The same compression of opportunity that makes them vulnerable to theft also makes them available as predatory labor. That dual role is central to the later emergence of guards.

\begin{corollary}[Comparative statics of patronage]
\label{cor:patronage-statics}
Fix $w_{it}$ and suppose $L_t = \phi_t \gamma \bar{w}^{H}_t$ and $o_{jt} = \lambda_t \gamma \bar{w}^{P}_t - c_T$. Whenever \eqref{eq:feasibility} holds, patron payoff \eqref{eq:patron-payoff} is increasing in $\phi_t$, decreasing in $\lambda_t$, decreasing in $k$, and decreasing in the attractiveness of personal theft only through the benchmark on the right-hand side of \eqref{eq:patron-dominance}. In particular, increases in delegation quality and decreases in poor outside options expand the set of states in which rich agents choose patronage over both self-theft and abstention.
\end{corollary}

\begin{proof}
See Appendix \ref{app:proofs}. The derivative of \eqref{eq:patron-payoff} with respect to $\phi_t$ is positive, the derivatives with respect to $\lambda_t$ and $k$ are negative, and the comparison against personal theft depends on whether these changes also widen the gap in \eqref{eq:patron-dominance}.
\end{proof}

These are exactly the margins varied later in the simulation figures. Patronage expands when delegated targeting becomes more effective, contracts when the outside option of poor workers strengthens, and weakens when contracting becomes expensive. The comparative statics in the code therefore match the comparative statics in the model.

\subsection{The switch from patronage to guarding}

If patronage succeeds for long enough, it changes the environment that made it profitable in the first place. Poor wealth is depleted, so $L_t$ declines over time as the supply of lucrative unguarded targets shrinks. At the same time, elite wealth rises, making concentrated fortunes worth defending. To capture the latter margin, let $\chi_t \in (0,1)$ denote the residual probability that elite wealth is targeted successfully when it is merely occupied but not formally guarded. This residual risk is low during active patronage because the patron remains at home and because dependent thieves have a reason not to antagonize the patrons who pay them. It need not remain low forever. Once patronage thins the poor prey base and contractual relations unravel, poaching against elites or opportunistic raids by former dependents become more attractive. Occupancy alone no longer replicates the discipline previously created by an active patronage network.

If elite agent $i$ hires a guard at wage $g$, and guarding blocks a raid with probability $q_t$, the expected value of protection is the avoided expected loss on concentrated wealth:
\begin{equation}
\pi^{\mathrm{guard}}_{it} = q_t \chi_t \gamma_E w_{it} - g,
\label{eq:guard-payoff}
\end{equation}
where $\gamma_E \in (0,1)$ is the fraction of elite wealth at stake in a successful elite raid. Poor workers accept guard employment provided the wage weakly exceeds their best alternative, which in the model means
\begin{equation}
g \ge \max\{o_{jt}, \underline{u}_{jt}\},
\label{eq:guard-IR}
\end{equation}
where $\underline{u}_{jt}$ is a noncriminal fallback option. The simulation imposes this acceptance condition directly when guards are hired from the poorest available agents.

The rich agent compares \eqref{eq:guard-payoff} with \eqref{eq:patron-payoff}. Guarding dominates patronage whenever the expected value of avoided elite loss exceeds the surplus from one more round of delegated predation:
\begin{equation}
q_t \chi_t \gamma_E w_{it} - g \ge \beta\left(L_t - k - o_{jt}\right).
\label{eq:guard-dominance}
\end{equation}

\begin{theorem}[Endogenous switch to guarding]
\label{thm:guarding}
Suppose guard labor is available at wage $g$ satisfying \eqref{eq:guard-IR} and patronage is initially feasible. Then guarding is optimal for elite agent $i$ whenever
\begin{equation}
w_{it} \ge \widehat{w}^{G}_t \equiv \frac{\beta\left(L_t - k - o_{jt}\right) + g}{q_t \chi_t \gamma_E}.
\label{eq:guard-threshold}
\end{equation}
Moreover, $\widehat{w}^{G}_t$ falls as the prey base among the poor is depleted, because $L_t - o_{jt}$ declines with $\bar{w}^{H}_t$ and $\bar{w}^{P}_t$.
\end{theorem}

\begin{proof}
See Appendix \ref{app:proofs}. Rearranging \eqref{eq:guard-dominance} yields the threshold. Depletion reduces the right-hand side of \eqref{eq:guard-dominance}, lowering the elite wealth level at which guarding becomes privately optimal.
\end{proof}

Several timing and role ambiguities also disappear under this threshold comparison. Rich households need not hire guards at the outset because elite wealth is not yet concentrated enough and active patronage already creates informal deterrence. Guards emerge later even though rich households were occupied earlier because occupancy no longer suffices once concentrated wealth becomes a tempting target and patronage ceases to discipline the population that once supplied predatory labor. The direction of predation also remains clear: during patronage, poor thieves are usually robbing poor households for rich principals, and once guarding dominates the same pool of poor agents may be reallocated from predation to protection.

\begin{corollary}[Comparative statics of guarding]
\label{cor:guard-statics}
The guarding threshold $\widehat{w}^{G}_t$ in \eqref{eq:guard-threshold} is decreasing in guard effectiveness $q_t$, decreasing in elite exposure $\chi_t$ and elite raid severity $\gamma_E$, increasing in the guard wage $g$, and increasing in the residual surplus from patronage $\beta(L_t-k-o_{jt})$. Guarding therefore becomes more likely when elite wealth is highly exposed and easily defended, and less likely when patronage remains highly profitable or guard labor becomes expensive.
\end{corollary}

\begin{proof}
See Appendix \ref{app:proofs}. The claims follow from direct differentiation of \eqref{eq:guard-threshold} with respect to the relevant parameters.
\end{proof}

Stronger guarding technology or greater elite exposure expands the region in which protection is privately optimal. Stronger patronage surplus or more expensive guard labor contracts that region. The move to policing is therefore an endogenous choice over technologies of control.

\subsection{Welfare}

The welfare criterion has to capture inequality and resource dissipation at the same time. An incomplete honest shock can move a society organized around reciprocal corruption toward more unequal and more resource-intensive regimes, while aggregate wealth alone remains uninformative because the benchmark conserves total wealth by construction. The analysis therefore uses discounted utilitarian welfare with explicit resource costs:
\begin{equation}
W_0 = \mathbb{E}_0 \sum_{t=0}^{\infty} \delta^t \mathcal{W}_t,
\qquad
\mathcal{W}_t = \sum_{i=1}^{N} u(w_{it}) - c_T n_t^{S} - k n_t^{H} - g n_t^{G} - \Psi_t,
\label{eq:welfare}
\end{equation}
where $u$ is strictly increasing and strictly concave, $n_t^{S}$ is the number of personal thieves, $n_t^{H}$ is the number of active patron-thief contracts, $n_t^{G}$ is the number of guards, and $\Psi_t$ denotes expected punishment losses.

\begin{proposition}[Redistribution, inequality, and resource cost]
\label{prop:welfare}
Suppose aggregate wealth $\sum_i w_{it}$ is fixed at date $t$ and $u$ is concave. Then any mean preserving spread of the wealth distribution weakly lowers $\sum_i u(w_{it})$. If, in addition, two regimes with the same aggregate wealth differ only in that one has weakly greater inequality and strictly greater values of $n_t^{S}$, $n_t^{H}$, $n_t^{G}$, or $\Psi_t$, then the latter regime yields strictly lower period welfare $\mathcal{W}_t$.
\end{proposition}

\begin{proof}
See Appendix \ref{app:proofs}. The first claim is Jensen's inequality. The second follows because the cost terms enter linearly and negatively in \eqref{eq:welfare}.
\end{proof}

A balanced corruption regime can already be normatively poor if theft effort and punishment risk consume resources, even when the distribution remains equal. Once the honest perturbation generates inequality and induces contracting and guarding costs, welfare can fall even if total wealth is approximately conserved in the short run. The relevant object is therefore distribution plus risk plus deadweight cost, and not redistribution alone.

\begin{corollary}[Guarding may improve welfare relative to patronage]
\label{cor:guard-welfare}
Fix the wealth distribution at date $t$. If adopting guarding lowers expected theft effort, contracting activity, and punishment losses by more than the resulting guard wage bill, then guarding raises $\mathcal{W}_t$ relative to patronage even though it need not restore the welfare of the original balanced benchmark.
\end{corollary}

\begin{proof}
See Appendix \ref{app:proofs}. The difference in period welfare is the reduction in theft-related costs minus the new guarding expenditure.
\end{proof}

Organized enforcement can be privately attractive for elites and even socially preferable to active patronage on a utilitarian measure, while still leaving society worse off than the initial equal benchmark once one accounts for the path by which elite wealth was accumulated. The model therefore leaves room for a limited welfare improvement from policing relative to patronage without mistaking that improvement for a return to egalitarian order.

A balanced theft cycle is a coherent benchmark in the unit transfer stage game. A single permanent honest perturbation breaks that benchmark in the repeated game when wealth-dependent exposure is strong enough. After symmetry breaks, proportional theft together with weak poor outside options makes patronage feasible and profitable. Later, depletion of the prey base and concentration of elite wealth make guarding privately optimal. The analytical results identify those margins, and the simulations trace their timing, prevalence, and welfare consequences.
